% 1.1
% Ficheiro incluído por main.tex

\section[Nome, rede e IP estático]{Nome, rede e IP estático}

\esquema{1-1-nome-rede-ip}{Windows Server: IP estático, gateway e DNS}

\passo{Passo 1 - Primeiro arranque}

\emph{Objetivo:} concluir o assistente inicial e dar o nome correto ao servidor.

Iniciar a VM e completar o assistente com os valores aconselhados:

\begin{itemize}
\item Região: \textbf{Portugal}; teclado: \textbf{Português}; idioma: \textbf{Inglês};
\item Palavra-passe do Administrator: \textbf{\texttt{\textless{}PASSWORD\_LAB\textgreater{}}} (a mesma em todo o laboratório).
\end{itemize}

\passo{Passo 2 - Mudar o nome do computador para \texttt{SR2-SRV}}

\textbf{Server Manager} → \textbf{Local Server} → clicar no nome atual do computador → \textbf{Change\ldots} → escrever \texttt{SR2-SRV} → \textbf{OK} → \textbf{Restart Now}.

Em alternativa, pelo PowerShell (como Administrator):

\begin{lstlisting}[style=psstyle]
Rename-Computer -NewName "SR2-SRV" -Restart
\end{lstlisting}

\emph{Verificação esperada:} após o reinício, barra de pesquisa do Windows → \textbf{About this computer} mostra o nome \texttt{SR2-SRV}.

\passo{Passo 3 - Abrir as definições do adaptador de rede}

\emph{Control Panel → Network and Internet → Network and Sharing Center → Change adapter settings} → botão direito no adaptador → \textbf{Properties} → \textbf{Internet Protocol Version 4 (TCP/IPv4)} → \textbf{Properties}.

\passo{Passo 4 - Atribuir o IP estático e testar a conectividade}

\emph{Objetivo:} dar ao DC um IP fixo e apontar o DNS para ele próprio.

\begin{center}\small
\begin{tabularx}{\linewidth}{|>{\raggedright\arraybackslash}p{0.32\linewidth}|>{\raggedright\arraybackslash}X|}
\hline
\rowcolor{azul!10}\textbf{Campo} & \textbf{Valor} \\
\hline
IP address & \texttt{\textless{}IP\_DC\textgreater{}} (192.168.X.Y) \\
\hline
Subnet mask & \texttt{255.255.255.0} \\
\hline
Default gateway & \texttt{192.168.X.1} \\
\hline
Preferred DNS & \texttt{\textless{}IP\_DC\textgreater{}} (o próprio servidor) \\
\hline
\end{tabularx}
\end{center}

\textbf{OK} para aplicar.

\begin{lstlisting}[style=psstyle]
ipconfig /all
ping 192.168.X.1
\end{lstlisting}

\emph{Verificação esperada:} \texttt{ipconfig} mostra \texttt{\textless{}IP\_DC\textgreater{}}, gateway \texttt{192.168.X.1} e DNS \texttt{\textless{}IP\_DC\textgreater{}}; o gateway responde ao ping.

\passo{Passo 5 - Aguardar o arranque dos serviços}

Aguardar que os serviços do servidor iniciem (o Server Manager abre sozinho) antes de começar a configuração.

\problemas{\hyperref[ec-a]{Erros Comuns - Parte A}}
